\documentclass[11pt,a4paper]{article}

% --- Packages ---
\usepackage[margin=.75in]{geometry}
\usepackage[utf8]{inputenc}
\usepackage[T1]{fontenc}
\usepackage{lmodern} % Fix for font not found error
\usepackage{titlesec}
\usepackage{enumitem}
\usepackage{hyperref}
\usepackage{xcolor}
\usepackage{fancyhdr}
\usepackage{comment}
\usepackage{cleveref}
\usepackage[bottom]{footmisc}
\usepackage{graphicx}
\usepackage{longtable}
\usepackage{booktabs}
\usepackage{array}
\usepackage{calc}
\usepackage{etoolbox}

% Standard Pandoc setup
\providecommand{\tightlist}{%
  \setlength{\itemsep}{0pt}\setlength{\parskip}{0pt}}

\crefformat{section}{\S#2#1#3}
\crefformat{subsection}{\S#2#1#3}
\crefformat{subsubsection}{\S#2#1#3}
\crefformat{figure}{#2Figure~#1#3}
\crefformat{footnote}{#2\footnotemark[#1]#3}

\linespread{0.9} % Riduce lo spazio tra le linee

% --- Colors ---
\definecolor{primary}{RGB}{0, 0, 0}
\definecolor{links}{HTML}{0000EE}
\definecolor{blue}{HTML}{26428B}

\hypersetup{
    colorlinks=true,
    linkcolor=links,
    urlcolor=links,
    citecolor=links,
}

% --- Section Formatting ---
\titleformat{\section}
  {\normalfont\Large\bfseries\color{blue}}{\thesection}{1em}{}
\titleformat{\subsection}
{\normalfont\large\bfseries\color{blue}}{\thesubsection}{1em}{}
\titleformat{\subsubsection}
{\normalfont\bfseries\color{blue}}{\thesubsubsection}{1em}{}

% --- Header and Footer ---
\pagestyle{fancy}
\fancyhead[L]{\small Giuseppe Capasso}
\fancyhead[R]{\small hashcat}
\fancyfoot[C]{\thepage}
\renewcommand{\headrulewidth}{0.4pt}

% --- Custom Title Command ---
\newcommand{\proposaltitle}[1]{
    \begin{center}
        \vspace*{0.1cm}
        {\LARGE \bfseries \color{blue} #1} \\[1.5ex]
        {\large \textbf{Giuseppe Capasso}} \\[1ex]
                \small{
            \vspace{0.1cm}
            \textbf{Materia:} hashcat\\
        }
            \end{center}
    \vspace{0.3cm}
    \hrule height 0.5pt
    \vspace{0.5cm}
}

\begin{document}

\proposaltitle{Password Cracking with Hashcat Lab}

\section{Password Cracking with Hashcat
Lab}\label{password-cracking-with-hashcat-lab}

\subsection{Introduction}\label{introduction}

In this lab, you will learn how to use Hashcat, a powerful password
recovery tool, to crack password hashes. Hashcat supports many hash
types and comes with various attack modes. This exercise will give you
practical experience with password cracking techniques in a controlled
environment.

\subsection{Objectives}\label{objectives}

\begin{itemize}
\tightlist
\item
  Understand how to use Hashcat to crack password hashes.
\item
  Learn about different hash types and Hashcat attack modes.
\end{itemize}

\subsection{Prerequisites}\label{prerequisites}

\begin{itemize}
\tightlist
\item
  Basic knowledge of command line interface (CLI) usage.
\item
  Access to a computer with Kali Linux installed.
\end{itemize}

\subsection{Required Tools}\label{required-tools}

\begin{itemize}
\tightlist
\item
  Kali Linux
\item
  Hashcat (pre-installed on Kali Linux)
\end{itemize}

\subsection{Lab Setup}\label{lab-setup}

\begin{enumerate}
\def\labelenumi{\arabic{enumi}.}
\item
  \textbf{Start your Kali Linux environment.} Ensure that you have
  administrative access.
\item
  \textbf{Open your terminal.} You will perform most of the tasks in
  this terminal.
\end{enumerate}

\subsection{Lab Tasks}\label{lab-tasks}

\subsubsection{Task 1: Identifying the Hash
Type}\label{task-1-identifying-the-hash-type}

Before you can begin cracking a hash, you need to identify what type of
hash it is. Use the \texttt{hash-identifier} tool that comes with Kali
Linux or refer to online resources to identify your hash type.

\textbf{Example Command:}

\begin{verbatim}
 hash-identifier 5f4dcc3b5aa765d61d8327deb882cf99
\end{verbatim}

\subsubsection{Task 2: Prepare a Password
List}\label{task-2-prepare-a-password-list}

For successful password cracking with Hashcat, you will need a list of
potential passwords, commonly referred to as a ``wordlist.'' Kali Linux
comes equipped with several wordlists, including the widely used
\texttt{rockyou.txt}.

\textbf{Steps to prepare your password list:}

\begin{enumerate}
\def\labelenumi{\arabic{enumi}.}
\tightlist
\item
  \textbf{Locate the Wordlist:} Kali Linux stores its wordlists in the
  \texttt{/usr/share/wordlists/} directory. For this task, we will use
  \texttt{rockyou.txt}, which is a comprehensive list of commonly used
  passwords.
\end{enumerate}

\begin{itemize}
\item
  \textbf{Command to locate} \texttt{rockyou.txt}\textbf{:}

\begin{verbatim}
ls /usr/share/wordlists/
\end{verbatim}
\end{itemize}

\begin{enumerate}
\def\labelenumi{\arabic{enumi}.}
\setcounter{enumi}{1}
\tightlist
\item
  Is the file there? You should see something like this:
\end{enumerate}

\begin{itemize}
\tightlist
\item
  ┌──(kali㉿kali)-{[}\textasciitilde{]} └─\$ ls /usr/share/wordlists -lh
  total 51M lrwxrwxrwx 1 root root 26 Mar 25 10:18 amass -\textgreater{}
  /usr/share/amass/wordlists lrwxrwxrwx 1 root root 25 Mar 25 10:18 dirb
  -\textgreater{} /usr/share/dirb/wordlists lrwxrwxrwx 1 root root 30
  Mar 25 10:18 dirbuster -\textgreater{} /usr/share/dirbuster/wordlists
  lrwxrwxrwx 1 root root 35 Mar 25 10:18 dnsmap.txt -\textgreater{}
  /usr/share/dnsmap/wordlist\_TLAs.txt lrwxrwxrwx 1 root root 41 Mar 25
  10:18 fasttrack.txt -\textgreater{}
  /usr/share/set/src/fasttrack/wordlist.txt lrwxrwxrwx 1 root root 45
  Mar 25 10:18 fern-wifi -\textgreater{}
  /usr/share/fern-wifi-cracker/extras/wordlists lrwxrwxrwx 1 root root
  28 Mar 25 10:18 john.lst -\textgreater{} /usr/share/john/password.lst
  lrwxrwxrwx 1 root root 27 Mar 25 10:18 legion -\textgreater{}
  /usr/share/legion/wordlists lrwxrwxrwx 1 root root 46 Mar 25 10:18
  metasploit -\textgreater{}
  /usr/share/metasploit-framework/data/wordlists lrwxrwxrwx 1 root root
  41 Mar 25 10:18 nmap.lst -\textgreater{}
  /usr/share/nmap/nselib/data/passwords.lst -rw-r--r-- 1 root root 51M
  May 12 2023 rockyou.txt.gz lrwxrwxrwx 1 root root 39 Mar 25 10:18
  sqlmap.txt -\textgreater{} /usr/share/sqlmap/data/txt/wordlist.txt
  lrwxrwxrwx 1 root root 25 Mar 25 10:18 wfuzz -\textgreater{}
  /usr/share/wfuzz/wordlist lrwxrwxrwx 1 root root 37 Mar 25 10:18
  wifite.txt -\textgreater{} /usr/share/dict/wordlist-probable.txt
\end{itemize}

\begin{enumerate}
\def\labelenumi{\arabic{enumi}.}
\setcounter{enumi}{2}
\tightlist
\item
  You need to decompress it:
\end{enumerate}

\begin{itemize}
\item
\begin{verbatim}
  sudo gzip -d /usr/share/wordlists/rockyou.txt.gz
      ```

  Note: If the file is already decompressed, this step can be skipped. You can check if it exists in uncompressed form by using the ls command.
  After decompressing the file, ensure it is ready for use by checking its content.
  ```bash
  head /usr/share/wordlists/rockyou.txt
\end{verbatim}
\end{itemize}

\subsubsection{Task 3: Cracking the
Hash}\label{task-3-cracking-the-hash}

Now that you have identified the type of hash and prepared a wordlist,
you're ready to use Hashcat to attempt to crack the hash. Follow the
steps below to properly configure and execute Hashcat.

\paragraph{Step 1: Prepare the Hash
File}\label{step-1-prepare-the-hash-file}

To begin, you'll need to place the hash you want to crack into a text
file. Make sure the hash format is compatible with the hash type you
identified in Task 1.

\textbf{Example Command to Create a Hash File:}

\begin{verbatim}
echo "5f4dcc3b5aa765d61d8327deb882cf99" > hash.txt
\end{verbatim}

\subsubsection{Step 2: Choose the Right Hashcat
Mode}\label{step-2-choose-the-right-hashcat-mode}

Identifying the correct mode is crucial for effectively using Hashcat to
crack a hash. Hashcat modes are specific to the type of hash you are
dealing with. The mode tells Hashcat how to properly interpret the hash
and what algorithm to use for cracking it.

\textbf{How to Choose the Mode:}

\begin{enumerate}
\def\labelenumi{\arabic{enumi}.}
\item
  \textbf{Refer to Hashcat's Mode List:} Hashcat documentation provides
  a complete list of modes for va
\item
  **Identify your hash type Recall the hash type you identified in Task
  1. Use this information to find the corresponding mode in the Hashcat
  list.
\end{enumerate}

\begin{verbatim}
hashcat --help | grep -i md5
\end{verbatim}

The output will show the mode number for MD5 hashes, typically 0. This
mode number is what you will use in the Hashcat command to crack the MD5
hash. Once you have identified the correct mode number from the list,
note it down as you will need to specify this mode in your Hashcat
command line to ensure that Hashcat processes the hash correctly.

\paragraph{Crack it!}\label{crack-it}

If you are cracking an MD5 has, you should use mode 0. Your hashcat
command will look like this:

\begin{verbatim}
hashcat -m 0 hash.txt /usr/share/wordlists/rockyou.txt
\end{verbatim}

\end{document}
